\section{The sole-progenitor feasibility theorem}\label{sec:theorem}

The strands of Sections~\ref{sec:endo}--\ref{sec:emp} assemble into the
paper's central result.

\begin{theorem}[Sole-progenitor feasibility]\label{thm:main}
There exists a parameter vector
\[
\theta^{*} \;=\; \bigl(T_0^{*},\ \gamma^{*},\ \Ne^{\max*},\ h^{*}\bigr)
\]
with $T_0^{*} \in [400, 595]$ kya, $\gamma^{*} \in [1.6, 2.0]$,
$\Ne^{\max*} \in [2\times10^{4}, 5\times10^{4}]$, and
$h^{*} \in [1.4, 1.9]\times10^{-3}$ --- the last interval bounded by the
heterozygosities of the observed human and chimpanzee genomes
\citep{PradoMartinez2013,Auton2015} --- such that the endowed-pair
trajectory $E(\theta^{*})$, which contains no pre-bottleneck population
and whose founding couple are the sole genetic progenitors of every
non-introgressed nucleotide of the descendant population, satisfies
simultaneously:
\begin{enumerate}
\item the pairwise-diversity identity $\pi = 2\mu E[T_2]$ at the observed
value ($\pi = 1.000\times10^{-3}$ vs.\ $10^{-3}$ at the canonical point,
Eq.~\eqref{eq:canonval});
\item the compression identity $\mathrm{TMR4A} = T_0$, matching the
direct-clock estimate $495 \pm 100$ kya \citep{Swamidass2019,Swamidass2022};
\item the observed autosomal coalescence-depth range
($0.5$--$2.0$ Myr typical, $\ge 2$ Myr at the deepest loci);
\item the strongly negative Tajima's $D$ of an expanding population
($D = -0.93$ at the canonical point, of the observed sign and order ---
the sign opposite to the one the classical bottleneck forecast predicted;
the residual to the observed series is addressed qualitatively,
Section~\ref{sec:limits});
\item the observed short-range linkage-disequilibrium profile
($E[r^2] \approx 0.40$--$0.30$ at $0.1$--$1$ kb, within the observed
$0.3$--$0.5$ range);
\item the uniparental identification of Proposition~\ref{prop:uni} with
published mtDNA/Y confidence intervals under the ancestral-rate
correction $s \ge 2.0$ of \citet{Scally2012} --- the theorem's one
clock-conditional condition (Section~\ref{sec:uni});
\item the invisibility of the bottleneck to PSMC/MSMC-type reconstruction
(Section~\ref{sec:ne});
\item the retention and regeneration of balanced-locus allelic richness
(Sections~\ref{sec:ret}--\ref{sec:poly});
\item under the descendant-clade extension, full nuclear coverage: the
introgressed archaic tracts are couple-descended through the archaic
branches, whose low diversity the model predicts (Theorem~\ref{thm:scope},
Figure~\ref{fig:branch}).
\end{enumerate}
with the archaic uniparental compartments scoped per
Theorem~\ref{thm:scope}.
\end{theorem}

\begin{proof}
Each numbered condition is established in the section cited: (1) by
Proposition~\ref{prop:master} and the canonical evaluation
\eqref{eq:canonval} with the three-way consistency check of
Section~\ref{sec:canonical}; (2) by Proposition~\ref{prop:tmr4a};
(3) by the depth decomposition of Section~\ref{sec:window}; (4)--(5) by
the simulations of Section~\ref{sec:sfsld}; (6) by Table~\ref{tab:clockscan};
(7) by Section~\ref{sec:ne}; (8) by Eqs.~\eqref{eq:hlaret}--\eqref{eq:hlares};
(9) by Theorem~\ref{thm:scope} and the branch sweep of
Section~\ref{sec:scope}.
Non-emptiness of the parameter region over the stated intervals is the
content of the $540$-point grid (Table~\ref{tab:gridextract},
Figure~\ref{fig:feasmap}).
\end{proof}

\subsection{Falsification conditions}
A feasibility theorem is only as good as its falsifiers, and the model's
are stated with the same precision as its claims:
\begin{enumerate}
\item a structure-robust revision of the direct-clock TMR4A estimate
outside $[395, 595]$ kya (the compression identity then forces
$T_0$ outside the diversity-feasible region, since feasibility requires
$T_0 \gtrsim 400$ kya at any admissible endowment);
\item resolution of the mutation-rate debate placing the ancestral-rate
correction at $s < 1.5$ (the corrected uniparental intervals would then
contain at most the window's youngest edge; Table~\ref{tab:clockscan}) ---
this removes the uniparental identification alone; the window's autosomal
support is untouched, Section~\ref{sec:uni};
\item demonstration that mid-Pleistocene $\Ne$ was persistently below
$2\times10^{4}$, with the growth-bias arguments against PSMC estimates
rebutted (the grid's $\Ne^{\max} = 10^{4}$ row is infeasible at every
$T_0$ in the window, at any admissible endowment);
\item a genome-wide excess of phase-aligned common-variant pairs above the
model's $50{:}50$ mixture prediction --- the endowment's two-point
signature of Section~\ref{sec:sfsld}, directly measurable in
phase-resolved data;
\item a robust direct-clock resolution of the archaic--modern population
split whose interval excludes $[550, 595]$ kya (scope S2 closes; scope S1
is untouched);
\item archaic genomes carrying heterozygosities well above the
small-branch band of Figure~\ref{fig:branch} (which would contradict the
descendant-clade reading of the introgressed fraction);
\item a conversion-corrected lineage count above the retention ceiling of
Eq.~\eqref{eq:hlaret} at any HLA-class locus (the per-locus test of
Section~\ref{sec:roots}), which falsifies the class at balanced loci
whatever the fate of the diversity-level arguments.
\end{enumerate}
Conditions 1--3 target the theorem's conditional layers (the TMR4A
reconstruction, the uniparental clock, the plateau size); conditions
4--7 target the model class itself. An adverse resolution of condition 1
or condition 3 empties the window and with it the theorem; condition 2
removes only the uniparental identification, leaving every direct-clock
condition intact (Section~\ref{sec:uni}) --- the asymmetry that makes the
uniparental layer corroborative rather than load-bearing. The theorem
fails outright if condition 1, 3, 4, or 7 holds; it narrows, without
failing, if conditions 2, 5, or 6 hold. Every condition is decidable on
public data by published methods --- an exposure the position opposite
the theorem does not share (Section~\ref{sec:asymmetry}).

\section{Objections answered}\label{sec:objections}

Table~\ref{tab:objections} runs the objection catalogue --- the standard
diversity and structure objections as posed in the modern debate
\citep{Gauger2012,Swamidass2019} and grounded in the classical bottleneck
literature \citep{Nei1975,Ayala1995} --- against the model's responses and
the sections that supply them.

\begin{table}[t]
\centering\small
\caption{The objection catalogue and the endowed-pair responses.}
\label{tab:objections}
\begin{tabular}{p{4.6cm}p{6.4cm}p{4.2cm}}
\toprule
Objection & Response & Section \\
\midrule
``$N=2$ alone carries $6\times$ too little diversity'' &
The figure assumes mutation-only diversity on a young bottleneck. With
the endowment, the depth reservoir is internal: $0.75h/2\mu \approx
1.23$ Myr $\ge$ the entire requirement &
\ref{sec:master}, \ref{sec:canonical} \\
``The couple must be sampled from
$\Nepre \approx 3\times10^4$'' &
That is the sampled-couple class. The EP class removes $\Nepre$ by
treating the couple's genomes as the initial condition; feasibility
holds without any pre-bottleneck population &
\ref{sec:classes}, \ref{sec:endo} \\
``The allele-frequency spectrum fits
only bottlenecks older than $\sim$500 kya'' &
Agreed and adopted: the joint window is $[400, 595]$ kya; the
diversity constraint no longer pulls the date younger &
\ref{sec:window} \\
``mtDNA and Y MRCAs are at 150--300
kya, not 500'' &
Under sole progenitorship both MRCAs \emph{are} the couple
(Prop.~\ref{prop:uni}); their published dates are anchored-clock
estimates requiring the published $s \approx 2$ ancestral-rate
correction; the uncorrected residual is the model's principal
conditional layer &
\ref{sec:uni} \\
``TMR4A $\approx 500$ kya contradicts a
recent couple'' &
The compression identity makes TMR4A $= T_0$ exactly; the estimate is
the model's own clock &
\ref{sec:window} \\
``PSMC/MSMC would show the bottleneck'' &
One-generation bottlenecks are invisible at reconstruction resolution;
the simulated profile shows a shallow depression, not $N=2$ &
\ref{sec:ne} \\
``Tajima's $D$ is too negative for a
bottlenecked population'' &
The EP model yields $D = -0.93$ at the canonical point --- the observed
sign and order, and the sign opposite to the classical bottleneck
forecast; the further deepening attributed to structure is argued from
the literature and runs toward the data (Sec.~\ref{sec:limits}) &
\ref{sec:sfsld} \\
``LD would overshoot under a
bottleneck-to-two'' &
The phase-mixed endowment gives $E[r^2_0] = 1/2$; the simulated
profile ($0.30$ at 1 kb) sits within the observed $0.3$--$0.5$ range ---
no overshoot is produced &
\ref{sec:sfsld} \\
``HLA lineages coalesce beyond the
human--chimp split'' &
Ayala's model priced a protracted bottleneck and ignored gene
conversion; one-generation retention is $2.1$--$2.5$ of 4 lineages;
regeneration supplies $8{,}705$ candidate lineages against
$\sim$$10^2$--$10^3$ observed; the deep-root lineage count remains the
open per-locus test (Sec.~\ref{sec:roots}) &
\ref{sec:poly} \\
``Archaic introgression violates sole
progenitorship'' &
The descendant-clade extension: the introgressing populations are
descendant branches of the couple, so admixture is gene flow within the
couple's pedigree; S2 (full nuclear coverage) holds at $T_0 \in [550,
595]$ kya, S1 (non-introgressed) unconditionally, and S3 (archaic
uniparentals included) is falsified clock-free &
\ref{sec:scope} \\
``Two individuals cannot be the sole
ancestors of Neanderthals, Denisovans, and modern humans alike'' &
Under the DCE the archaic clades branch from the couple's descendants;
their low diversity, deep cross-coalescence with moderns, and admixture
date all follow from the same founding event &
\ref{sec:scope} \\
\bottomrule
\end{tabular}
\end{table}

\section{The nature of the proof}\label{sec:nature}
The theorem proved here is a claim of
\emph{logical compatibility}, not of actuality: the diversity and
structure of the human genome do not falsify a sole-progenitor couple in
the window $[400, 595]$ kya. The distinction matters for what the
argument can and cannot do. It cannot show that the couple existed; it
can show --- and does --- that the population-genetic objections most
often cited against the claim are either incorrect as stated, or correct
only for a different model class, or correct as conditional layers that
the model states openly.

The proof's structure also clarifies the relationship between the two
roads the debate has taken. The sampled-couple class proves a
bottleneck-survival theorem: a couple \emph{carries} the diversity of a
pre-existing population across a
one-generation bottleneck. The present paper proves the feasibility
theorem in the endowed-pair class: a couple \emph{sources} the diversity,
if its genomes carry a divergence within the range that primate genomes
in fact exhibit. The ``$\Nepre \ge 3\times10^{4}$'' that forced the
carrier reading is not a fact about couples; it is a fact about
\emph{randomly sampled} couples --- an artifact of the sampling
assumption. Once the couple's genomes enter as an initial condition bounded
by biology rather than by ancestry, the same equations that demanded a
pre-bottleneck population are satisfied without one, at a bottleneck time
that three independent autosomal arguments already preferred. Both model
classes remain on the table; the diversity data as they stand do not
choose between them, because both generate the same observables within
their windows. What the data \emph{do} constrain --- the bottleneck time,
the plateau size, the endowment magnitude, the scoping of the archaic
fraction --- are stated quantitatively.

The provenance of the endowment is the one question the mathematics
cannot touch, and the paper's claim about it is underdetermination, not
resolution. Two founders whose genomes carry $\sim$4.9 million SNVs
each --- the coalescent-equivalent of $\Ne \approx 3.3\times10^{4}$
(Section~\ref{sec:bound}) --- may have acquired that divergence in three
ways: they may be \emph{sampled} from a population large enough to
generate it (the sampled-couple class); they may be \emph{constituted de
novo} at that divergence; or their genomes may be \emph{assembled} from
diverged ancestral lineages, the way a hybrid's two haplotypes carry the
divergence of their source populations. The genetics of their
descendants is identical in the three cases --- that is the content of
Eqs.~\eqref{eq:sampled} and \eqref{eq:master} sharing the same form ---
so no diversity measurement can adjudicate the provenance, and a theory
that pretends to is importing an assumption, not deriving a result. The
paper's contribution is to have shown that the quantitative case against
sole progenitorship, as posed in the diversity literature, does not exist
at the stated strength; which provenance reading a reader accepts is not
a population-genetic question.

\subsection{The falsifiability asymmetry}\label{sec:asymmetry}
The theorem's scientific standing is easiest to see from the shape of its
exposure. Every constraint it satisfies is a mainstream result read on
the direct clock; every correction it applies is a published
reconciliation; its one free parameter is bounded by the heterozygosities
of genomes that exist. Its falsifiers, accordingly, are measurable
quantities, stated in advance and decidable on public data
(Section~\ref{sec:theorem}): a revised TMR4A estimate, a resolved rate
debate, a demonstrated plateau ceiling, a phase-resolved excess, an
archaic--modern split interval, an archaic heterozygosity level, a
corrected lineage count. This is the exposure that origin claims are
usually said to avoid, and the model accepts it deliberately: seven
numbers, every one of them measurable today, stand between it and
refutation.

The position opposite the theorem does not face the same exposure. To
maintain the diversity objection against the feasibility result, a critic
must conjoin at least four assumptions, and the conjunction is the point.
(i) That the founding couple's genomes are a random sample of a
pre-existing population. The identical-observables result --- Eqs.~\eqref{eq:sampled}
and \eqref{eq:master} sharing one form --- makes that assumption
undecidable by any diversity measurement; it is not established by the
data, and it is not refutable by the data, which is the definition of an
assumption carried rather than derived. It is also, verbatim, the
assumption sole progenitorship denies, so the objection rests on assuming
what is at issue. (ii) That the mid-Pleistocene plateau was persistently
of order $10^{4}$. This requires discounting the transient peaks MSMC
infers \citep{LiDurbin2011,SchiffelsDurbin2014}, the growth-bias
arguments against low ancient-plateau readings, and the independent
severe-bottleneck inference of \citet{Hu2023} --- and reconstructions
that disagree with one another by factors of several cannot be read
selectively without an argument the literature does not currently supply.
(iii) That the ancestral-rate correction $s \approx 2$, documented on the
nuclear genome and required to reconcile pedigree with phylogeny at all
\citep{Scally2012}, does not transfer to the uniparental clocks. No
published measurement establishes this, and none currently could: there
is no pedigree calibration of archaic uniparental rates. (iv) That
HLA-class sequence depths are clean coalescent signal. The conversion
literature itself discredits that default \citep{GarriganHedrick2003},
and the protocol that would settle the question has never been run on
either side (Section~\ref{sec:roots}).

No single assumption in this list is unreasonable to entertain; the
objection's dependence on their conjunction is what weakens it. Assumption
(i) is untestable by the data the debate is conducted on; (ii) is
selective among reconstructions that disagree among themselves; (iii) is
untestable on present data; (iv) is untested where a test exists. The
falsifiability gradient therefore runs against the objection, not the
theorem: the theorem names the numbers that would destroy it, while the
objection holds only if a conjunction of assumptions --- none
individually established, one in principle unestablishable on diversity
data --- is granted silently. The paper's claim to scientific consistency
is exactly this: on every quantity the data can measure --- $\pi$,
coalescence depths, spectral sign, linkage disequilibrium, TMR4A, the
plateau band, lineage counts --- the sole-progenitor claim in its stated
scope is unrebutted, and the remaining disagreement sits on assumptions
the data cannot reach.

\section{Scope and limitations}\label{sec:limits}
Six limitations bound the theorem.
\begin{enumerate}
\item \emph{The uniparental conditional layer.} The model's principal
tension is the anchored-clock dating of the mtDNA and Y MRCAs
(Section~\ref{sec:uni}). The published ancestral-rate correction
\citep{Scally2012} closes it at $s \ge 2.0$ --- a factor calibrated on
nuclear divergence, whose transfer to the single-locus uniparental clocks
is itself unproven and is part of the layer's conditionality
(Section~\ref{sec:uni}); if the correction is
rejected, the uniparental dates prefer $T_0 \approx 200$--$300$ kya,
which the EP class admits only at $\Ne^{\max} \ge 2\times10^{4}$ with the
upper endowment ($\pi \approx 0.9$--$1.0\times10^{-3}$) --- feasible for
diversity but outside the TMR4A window unless the compression estimate is
itself structure-inflated by $\sim$1.5--$2\times$ in the direction
documented for structure-biased coalescent inference
\citep{ErikssonManica2012,Scerri2019}. The two corrections trade off; no
parameterization known to us closes both at once without one of them,
and the paper does not claim otherwise.
\item \emph{The plateau size.} Feasibility requires
$\Ne^{\max} \ge 2\times10^{4}$ --- the model's second conditional layer,
and a strong one. Published long-term reconstructions for human ancestry
frequently sit near $\Ne \approx 10^{4}$, a plateau at which the grid is
infeasible at every $T_0$ (Table~\ref{tab:gridfull}); the requirement is
met only at the upper end of the published reconstructions --- the
transient peaks MSMC infers \citep{LiDurbin2011,SchiffelsDurbin2014} and
the growth-bias arguments that discount low ancient-plateau readings ---
and if the low-end readings are taken at face value, the EP class fails.
The theorem is therefore conditional on that reading of the demographic
literature, stated with its falsifier (condition 3,
Section~\ref{sec:theorem}).
\item \emph{The endowment provenance.} The model takes $h$ as an initial
condition (Section~\ref{sec:nature}); its generation mechanism is outside
population genetics.
\item \emph{Backbone approximation.} The backbone's segregating fraction
is computed under a four-allele approximation; mixed 2/3-allele
configurations shift $\theta_W$ by a few percent and $D$ by less.
\item \emph{Structure not simulated jointly.} The structure-dependent
refinements (deeper $D$, spectrum corrections, TMR4A inference bias) are
argued from the cited literature rather than simulated in this paper's
grid; the Tajima's-$D$ claim in particular is open at exactly this point
--- if a structured simulation of the EP class fails to deepen $D$
materially below the canonical $-0.93$, the spectrum condition holds in
sign and order only. A structured joint simulation is the natural next
step.
\item \emph{Scoping and the descendant-clade layer.} The theorem covers
the non-introgressed genome unconditionally (scope S1); full nuclear
coverage (S2, the descendant-clade extension) is conditional on
$T_0 \in [550, 595]$ kya and on archaic uniparentals replaced from
outside the couple's descent --- which makes sole progenitorship a claim
about modern genomes, not about the archaic clades' uniparental
compartments (Section~\ref{sec:scope}) --- and scope S3 is falsified
(Theorem~\ref{thm:scope}). The $\sim$2--6\%
introgressed fraction is a shared boundary of every model in the debate.
The DCE's split condition sits at the youngest edge of the published
archaic--modern divergence estimates ($550$--$765$ kya,
\citealp{Pruefer2014}); a robust resolution excluding that edge narrows
the claim to S1 without touching it.
\end{enumerate}

These six prices are the theorem's own, and they are quantified, bounded,
and carried with falsifiers. Their mirror image belongs in the same
accounting: each has a symmetric assumption on the objection's side ---
the sampled reading no diversity measurement can test
(Section~\ref{sec:nature}), the persistently low plateau that must
discount the upper reconstructions and the growth-bias arguments, the
refusal of a published correction (limit 1), the uncorrected HLA depths
(Section~\ref{sec:roots}). A reader weighing the two positions weighs a
priced conditionality against an unpriced one
(Section~\ref{sec:asymmetry}).

\section{Conclusion}\label{sec:conclusion}
The paper's central claim, in one sentence: \emph{a founding couple of
$N=2$ individuals, carrying a genomic endowment within the range that
natural primate genomes exhibit, expanding rapidly at $T_0 \approx 550$
kya, reproduces the observed diversity, frequency spectrum,
coalescence-depth distribution, linkage-disequilibrium profile,
uniparental genealogy structure, and balanced-locus richness of modern
humans --- and, under the descendant-clade extension, accounts for the
introgressed archaic fraction as descendant-branch gene flow, with the
archaic branches' low diversity predicted rather than postulated --- with
no pre-bottleneck population in the model class: the endowment does
numerically what the sampled-couple class does with its $\Nepre$, the two
classes generate identical observables, and the data as they stand do not
choose between them. What the theorem establishes is underdetermination
under exact bounds: the diversity objection fails as stated, the
admissible window is $[400, 595]$ kya, at least $47\%$ of today's
diversity must be carried into the founding event ($h \ge
0.63\times10^{-3}$), and the endowment's provenance --- sampling, de-novo
constitution, or assembly from diverged lineages --- lies beyond any
diversity measurement, which is why the historical objection never
followed from the data.} The
argument required no new genetics: it required the removal of one
assumption --- that the couple is a random sample --- and the honest
pricing of the two conditional layers that remain: the uniparental clock
calibration, a conditional corroboration that the window's autosomal
evidence does not depend on, and the plateau size, a strong requirement
met only at the upper end of the published reconstructions --- stated
with their falsifiers. The
historical diversity objection, under which two individuals were said to
carry too little variation to found a species, is thereby answered in its
strongest form: the diversity two primates can carry is not a shortfall
against the requirement but a reservoir exceeding it, and the structure
of the modern genome --- spectrum, depths, disequilibria, balanced loci,
archaic branches --- is what that reservoir plus half a million years of
mutation and drift predicts. The position that would restore the
objection must grant, silently, what the data cannot test; the position
defended here asks only what the data can measure. On the measurable
quantities the theorem stands: a sole-progenitor couple at
$T_0 \approx 550$ kya is consistent with every mainstream constraint the
human genome supplies, its conditional layers are priced and carry
falsifiers, and the historical diversity objection --- the reason the
claim was dismissed --- does not survive the removal of its one hidden
assumption.

\section*{Data and code availability}

Every quantitative result in this paper --- the master-equation
evaluation, the $540$-point feasibility grid, the msprime coalescent
simulations, the site-frequency spectrum and Tajima's $D$, the
linkage-disequilibrium decay, the effective-size profile, the HLA
regeneration calculus, the scoping tests, and the empirical anchors --- is
emitted by a single fixed-seed Python script,
\texttt{sole\_progenitor\_model.py}, which executes end to end in
${\approx}\,3$ minutes on one core (\texttt{python3
sole\_progenitor\_model.py}; the \texttt{--quick} flag gives a
${\approx}\,20$-second check). The complete package --- the script, all
simulation outputs, the pinned requirements file, the citation-audit
report, and the full \LaTeX{} source of the manuscript --- is publicly
available at
\url{https://github.com/Atikur-rahman-ar/sole-genetic-progenitor}. All
stochastic components use the fixed seeds listed in
Appendix~\ref{app:repro}, so every number in this paper is
bit-reproducible from the repository alone. The paper generates no new
sequence data and reads no proprietary or restricted data: every empirical
quantity it uses is taken from the published studies cited in the text,
and the corresponding values are tabulated as explicit anchors in
\texttt{outputs/empirical\_anchors.csv} of the repository.

\begin{thebibliography}{60}
\setlength{\itemsep}{1pt}

\bibitem[Auton et al.(2015)]{Auton2015}
Auton, A., et al. (2015). A global reference for human genetic variation.
\emph{Nature} 526:68--74.

\bibitem[Ayala(1995)]{Ayala1995}
Ayala, F.J. (1995). The myth of Eve: molecular biology and human origins.
\emph{Science} 270:1930--1936.

\bibitem[Barton and Charlesworth(1984)]{BartonCharlesworth1984}
Barton, N.H. \& Charlesworth, B. (1984). Genetic revolutions, founder
effects, and speciation. \emph{Annual Review of Ecology and Systematics}
15:133--164.

\bibitem[Cann et al.(1987)]{Cann1987}
Cann, R.L., Stoneking, M. \& Wilson, A.C. (1987). Mitochondrial DNA and
human evolution. \emph{Nature} 325:31--36.

\bibitem[Crow and Kimura(1970)]{CrowKimura1970}
Crow, J.F. \& Kimura, M. (1970). \emph{An Introduction to Population
Genetics Theory}. Harper \& Row, New York.

\bibitem[Eriksson and Manica(2012)]{ErikssonManica2012}
Eriksson, A. \& Manica, A. (2012). Effect of ancient population structure
on the degree of polymorphism shared between modern human populations and
ancient hominins. \emph{Proceedings of the National Academy of Sciences
USA} 109, doi:10.1073/pnas.1200567109.

\bibitem[Ewing et al.(2008)]{Ewing2008}
Ewing, S.R., et al. (2008). Inbreeding and loss of genetic variation in a
reintroduced population of Mauritius kestrel. \emph{Conservation Biology}
22:395--404.

\bibitem[Fenner(2005)]{Fenner2005}
Fenner, J.N. (2005). Cross-cultural estimation of the human generation
interval for use in genetics-based population divergence studies.
\emph{American Journal of Physical Anthropology} 128:415--423.

\bibitem[Garrigan and Hammer(2006)]{GarriganHammer2005}
Garrigan, D. \& Hammer, M.F. (2006). Reconstructing human origins in the
genomic era. \emph{Nature Reviews Genetics} 7:669--680.

\bibitem[Garrigan and Hedrick(2003)]{GarriganHedrick2003}
Garrigan, D. \& Hedrick, P.W. (2003). Detecting adaptive molecular
polymorphism: lessons from the MHC. \emph{Evolution} 57:1707--1712.

\bibitem[Gauger et al.(2012)]{Gauger2012}
Gauger, A.K., Axe, D. \& Luskin, C. (2012). \emph{Science and Human Origins}.
Discovery Institute Press, Seattle.

\bibitem[Green et al.(2008)]{Green2008}
Green, R.E., et al. (2008). A complete Neandertal mitochondrial genome
sequence determined by high-throughput sequencing. \emph{Cell}
134:416--426.

\bibitem[Hill and Weir(1988)]{HillWeir1988}
Hill, W.G. \& Weir, B.S. (1988). Variances and covariances of squared
linkage disequilibria for finite populations. \emph{Theoretical
Population Biology} 33:54--78.

\bibitem[Hoelzel et al.(2024)]{Hoelzel2024}
Hoelzel, A.R., et al. (2024). Genomics of post-bottleneck recovery in the
northern elephant seal. \emph{Nature Ecology \& Evolution} 8:686--694.

\bibitem[H\"ossjer and Gauger(2019)]{HossjerGauger2019}
H\"ossjer, O. \& Gauger, A.K. (2019). A single-couple human origin is
possible. \emph{BIO-Complexity} 2019(1):1--20. doi:10.5048/BIO-C.2019.1.

\bibitem[Hu et al.(2023)]{Hu2023}
Hu, W., Hao, Z., Du, P., Di, F.Z., et al. (2023). Genomic inference of a
severe human bottleneck during the Early to Middle Pleistocene transition.
\emph{Science} 381:979--984.

\bibitem[Hudson(1990)]{Hudson1990}
Hudson, R.R. (1990). Gene genealogies and the coalescent process.
\emph{Oxford Surveys in Evolutionary Biology} 7:1--44.

\bibitem[Jackson et al.(2022)]{Jackson2022}
Jackson, H.A., et al. (2022). Genomic erosion in a demographically
recovered bird species during conservation rescue. \emph{Conservation
Biology} 36:e13918.

\bibitem[J\'onsson et al.(2017)]{Jonsson2017}
J\'onsson, H., et al. (2017). Parental influence on human germline de novo
mutations in 1,548 trios from Iceland. \emph{Nature} 549:519--522.

\bibitem[Kaeuffer et al.(2007)]{Kaeuffer2007}
Kaeuffer, R., Coltman, D.W., Chapuis, J.-L., Pontier, D. \& R\'eale, D.
(2007). Unexpected heterozygosity in an island mouflon population founded
by a single pair of individuals. \emph{Proceedings of the Royal Society B}
274:527--533.

\bibitem[Karmin et al.(2016)]{Karmin2016}
Karmin, M., et al. (2016). A recent bottleneck of Y chromosome diversity
coincides with a global change in culture. \emph{Genome Research}
26:459--466.

\bibitem[Kong et al.(2012)]{Kong2012}
Kong, A., et al. (2012). Rate of de novo mutations and the importance of
father's age to disease risk. \emph{Nature} 488:471--475.

\bibitem[Li and Durbin(2011)]{LiDurbin2011}
Li, H. \& Durbin, R. (2011). Inference of human population history from
individual whole-genome sequences. \emph{Nature} 475:493--496.

\bibitem[McMullan et al.(2018)]{McMullan2018}
McMullan, M., et al. (2018). The ash dieback invasion of Europe was
founded by two genetically divergent individuals. \emph{Nature Ecology \&
Evolution} 2:1000--1008.

\bibitem[Mendez et al.(2013)]{Mendez2013}
Mendez, F.L., et al. (2013). An African American paternal lineage adds an
extremely ancient root to the human Y chromosome phylogenetic tree.
\emph{American Journal of Human Genetics} 92:454--459.

\bibitem[Mendez et al.(2016)]{Mendez2016}
Mendez, F.L., Poznik, G.D., Castellano, S. \& Bustamante, C.D. (2016).
The divergence of Neandertal and modern human Y chromosomes.
\emph{American Journal of Human Genetics} 98:728--734.

\bibitem[Meyer et al.(2012)]{Meyer2012}
Meyer, M., et al. (2012). A high-coverage genome sequence from an archaic
Denisovan individual. \emph{Science} 338:222--226.

\bibitem[Nei et al.(1975)]{Nei1975}
Nei, M., Maruyama, T. \& Chakraborty, R. (1975). The bottleneck effect and
genetic variability in populations. \emph{Evolution} 29:1--10.

\bibitem[Posth et al.(2017)]{Posth2017}
Posth, C., Wi\ss{}ing, C., Kitagawa, K., et al. (2017). Deeply divergent
archaic mitochondrial genome provides lower time boundary for African gene
flow into Neanderthals. \emph{Nature Communications} 8:16046.

\bibitem[Poznik et al.(2013)]{Poznik2013}
Poznik, G.D., et al. (2013). Sequencing Y chromosomes resolves discrepancy
in time to common ancestor of males versus females. \emph{Science}
341:562--565.

\bibitem[Prado-Martinez et al.(2013)]{PradoMartinez2013}
Prado-Martinez, J., et al. (2013). Great ape genetic diversity and
population history. \emph{Nature} 499:471--475.

\bibitem[Pr\"ufer et al.(2014)]{Pruefer2014}
Pr\"ufer, K., et al. (2014). The complete genome sequence of a Neanderthal
from the Altai Mountains. \emph{Nature} 505:43--49.

\bibitem[Rogers et al.(2020)]{Rogers2020}
Rogers, A.R., Harris, N.S. \& Achenbach, J.D. (2020). Neanderthal-Denisovan
ancestors interbred with a distantly related hominin. \emph{Science
Advances} 6:eaay5483.

\bibitem[Sankararaman et al.(2012)]{Sankararaman2012}
Sankararaman, S., et al. (2012). The date of interbreeding between
Neandertals and modern humans. \emph{PLoS Genetics} 8:e1002947.

\bibitem[Sankararaman et al.(2014)]{Sankararaman2014}
Sankararaman, S., et al. (2014). The genomic landscape of Neanderthal
ancestry in present-day humans. \emph{Nature} 507:354--357.

\bibitem[Scally and Durbin(2012)]{Scally2012}
Scally, A. \& Durbin, R. (2012). Revising the human mutation rate:
implications for understanding human evolution. \emph{Nature Reviews
Genetics} 13:745--753.

\bibitem[Scerri et al.(2019)]{Scerri2019}
Scerri, E.M.L., Chikhi, L. \& Thomas, M.G. (2019). Beyond multiregional
and simple out-of-Africa models of human evolution. \emph{Nature Ecology
\& Evolution} 3:1370--1372.

\bibitem[Schiffels and Durbin(2014)]{SchiffelsDurbin2014}
Schiffels, S. \& Durbin, R. (2014). Inferring human population size and
separation history from multiple genome sequences. \emph{Nature Genetics}
46:919--925.

\bibitem[Slatkin and Excoffier(2012)]{SlatkinExcoffier2012}
Slatkin, M. \& Excoffier, L. (2012). Serial founder effects during range
expansion: a spatial analog of genetic drift. \emph{Genetics}
191:171--181.

\bibitem[Swamidass(2019)]{Swamidass2019}
Swamidass, S.J. (2019). \emph{The Genealogical Adam and Eve: The Surprising
Science of Universal Ancestry}. IVP Academic, Downers Grove, IL.

\bibitem[Swamidass(2022)]{Swamidass2022}
Swamidass, S.J. (2022). The misunderstood science of genetic bottlenecks.
Paper presented at the 2022 Annual Meeting; \emph{Peaceful Science}.

\bibitem[Sz\'ucs et al.(2017)]{Szucs2017}
Sz\'ucs, M., Melbourne, B.A., Tuff, T., Weiss-Lehman, C. \& Hufbauer,
R.A. (2017). Genetic and demographic founder effects have long-term
fitness consequences for colonising populations. \emph{Ecology Letters}
20:436--444.

\bibitem[Tajima(1989)]{Tajima1989}
Tajima, F. (1989). Statistical method for testing the neutral mutation
hypothesis by DNA polymorphism. \emph{Genetics} 123:585--595.

\bibitem[Takahata(1990)]{Takahata1990}
Takahata, N. (1990). A simple genealogical structure of strongly balanced
allelic lines and trans-species evolution of polymorphism.
\emph{Proceedings of the National Academy of Sciences USA} 87:2419--2423.

\bibitem[Tenesa et al.(2007)]{Tenesa2007}
Tenesa, A., et al. (2007). Recent human effective population size
estimated from linkage disequilibrium. \emph{American Journal of Human
Genetics} 80:495--505.

\bibitem[Tennessen et al.(2012)]{Tennessen2012}
Tennessen, J.A., et al. (2012). Evolutionary history and adaptation from
high-coverage whole-genome sequence data. \emph{Science} 337:64--74.

\bibitem[Vigilant et al.(1991)]{Vigilant1991}
Vigilant, L., Stoneking, M., Harpending, H., Hawkes, K. \& Wilson, A.C.
(1991). African populations and the evolution of human mitochondrial DNA.
\emph{Science} 253:1503--1507.

\bibitem[von Salom\'e et al.(2007)]{vonSalome2007}
von Salom\'e, J., Gyllensten, U. \& Bergstr\"om, T.F. (2007). Full-length
sequence analysis of the HLA-DRB1 locus suggests a recent origin of
alleles. \emph{Immunogenetics} 59:261--271.

\bibitem[von Seth et al.(2022)]{vonSeth2022}
von Seth, J., et al. (2022). Genomic trajectories of a near-extinction
event in the Chatham Island black robin. \emph{BMC Genomics} 23:747.

\bibitem[Watterson(1975)]{Watterson1975}
Watterson, G.A. (1975). On the number of segregating sites in genetical
models without recombination. \emph{Theoretical Population Biology}
7:256--276.

\bibitem[Wright(1931)]{Wright1931}
Wright, S. (1931). Evolution in Mendelian populations. \emph{Genetics}
16:97--159.

\bibitem[Xue et al.(2018)]{Xue2018}
Xue, D., Graves, J., Carranza, A., et al. (2018). Successful worldwide
invasion of the veined rapa whelk, \emph{Rapana venosa}, despite a
dramatic genetic bottleneck. \emph{Biological Invasions}
20:3297--3314.

\end{thebibliography}
